\chapter{Prefect Workflow Orchestration}
\label{ch:prefect}

\section{Overview}

Prefect 2.0 is used to orchestrate the entire ML pipeline, providing task scheduling, error handling, retries, and comprehensive logging.

\section{Prefect Architecture}

\subsection{Core Concepts}

\begin{itemize}
    \item \textbf{Flow:} A collection of tasks and their relationships
    \item \textbf{Task:} An individual unit of work
    \item \textbf{Run:} A single execution of a flow
    \item \textbf{Deployment:} A flow packaged with schedule and config
    \item \textbf{Work Queue:} A queue for distributing work to agents
\end{itemize}

\subsection{Task Dependencies Graph}

\begin{figure}[H]
\centering
\begin{tikzpicture}[scale=0.9, node distance=2.2cm]
    \node[draw, fill=blue!30, rounded rectangle, minimum width=2cm] (1) at (3, 10) {Start};
    
    \node[draw, fill=green!30, minimum width=2cm, minimum height=0.7cm] (2) at (3, 8.5) {Fetch Data};
    \node[draw, fill=green!30, minimum width=2cm, minimum height=0.7cm] (3) at (3, 7) {Validate Data};
    
    \node[draw, fill=yellow!30, minimum width=2cm, minimum height=0.7cm] (4) at (1, 5) {Engineer};
    \node[draw, fill=yellow!30, minimum width=2cm, minimum height=0.7cm] (5) at (3, 5) {Preprocess};
    \node[draw, fill=yellow!30, minimum width=2cm, minimum height=0.7cm] (6) at (5, 5) {Split Data};
    
    \node[draw, fill=orange!30, minimum width=2cm, minimum height=0.7cm] (7) at (1, 3) {Train Clf};
    \node[draw, fill=orange!30, minimum width=2cm, minimum height=0.7cm] (8) at (3, 3) {Train Reg};
    \node[draw, fill=orange!30, minimum width=2cm, minimum height=0.7cm] (9) at (5, 3) {Train DL};
    
    \node[draw, fill=purple!30, minimum width=2cm, minimum height=0.7cm] (10) at (2, 1) {Evaluate};
    \node[draw, fill=purple!30, minimum width=2cm, minimum height=0.7cm] (11) at (4, 1) {Version};
    
    \node[draw, fill=red!30, rounded rectangle, minimum width=2cm] (12) at (3, -1) {End};
    
    \draw[->] (1) -- (2);
    \draw[->] (2) -- (3);
    \draw[->] (3) -- (4);
    \draw[->] (3) -- (5);
    \draw[->] (3) -- (6);
    \draw[->] (4) -- (7);
    \draw[->] (5) -- (8);
    \draw[->] (6) -- (9);
    \draw[->] (7) -- (10);
    \draw[->] (8) -- (10);
    \draw[->] (9) -- (10);
    \draw[->] (10) -- (11);
    \draw[->] (11) -- (12);
\end{tikzpicture}
\caption{Prefect Task Dependency Graph}
\label{fig:prefect_dag}
\end{figure}

\section{Main ML Pipeline Flow}

\subsection{Flow Definition}

\begin{lstlisting}[language=Python, caption={Prefect ML Pipeline Definition (Simplified)}]
from prefect import flow, task
from prefect.tasks.prefect_orm import get_run_context
import pandas as pd
from src.fetch_bitcoin_data import fetch_bitcoin_data
from src.preprocess_bitcoin import preprocess_bitcoin_data
from src.train_all_models import train_all_models

@task(retries=3, retry_delay_seconds=30)
def fetch_data_task(days: int = 365) -> pd.DataFrame:
    """Fetch Bitcoin data from CoinGecko API"""
    print(f"Fetching {days} days of Bitcoin data...")
    df = fetch_bitcoin_data(days=days)
    print(f"✓ Fetched {len(df)} records")
    return df

@task(retries=2)
def validate_data_task(df: pd.DataFrame) -> bool:
    """Validate data quality"""
    assert len(df) > 0, "Empty dataframe"
    assert df['price'].notna().sum() > 0, "No prices"
    print(f"✓ Data validation passed ({len(df)} records)")
    return True

@task
def engineer_features_task(df: pd.DataFrame) -> pd.DataFrame:
    """Engineer 24 technical indicators"""
    print("Engineering features...")
    df_features = engineer_features(df)
    print(f"✓ Engineered {df_features.shape[1]} features")
    return df_features

@task
def preprocess_task(df: pd.DataFrame) -> tuple:
    """Preprocess data"""
    print("Preprocessing...")
    df_processed, scaler = preprocess_bitcoin_data(df)
    print(f"✓ Preprocessing complete: {df_processed.shape}")
    return df_processed, scaler

@task
def split_data_task(df: pd.DataFrame, test_size: float = 0.2) -> tuple:
    """Temporal train/test split"""
    split_idx = int(len(df) * (1 - test_size))
    X_train = df.iloc[:split_idx]
    X_test = df.iloc[split_idx:]
    print(f"✓ Split: {len(X_train)} train, {len(X_test)} test")
    return X_train, X_test

@task(retries=2, retry_delay_seconds=60)
def train_models_task(X_train, X_test, y_train_clf, y_train_reg) -> dict:
    """Train all models (classification, regression, deep learning)"""
    print("Training models...")
    
    results = train_all_models(
        X_train, X_test,
        y_train_clf, y_train_reg,
        epochs=100, batch_size=32
    )
    
    print(f"✓ Trained {len(results)} models")
    return results

@task
def evaluate_models_task(results: dict) -> dict:
    """Evaluate and rank models"""
    print("Evaluating models...")
    
    best_model = max(results.items(), 
                    key=lambda x: x[1]['metrics']['accuracy'])
    
    print(f"✓ Best model: {best_model[0]} ({best_model[1]['metrics']['accuracy']:.2%})")
    return results

@task
def version_models_task(results: dict, models: dict) -> str:
    """Version and save models to registry"""
    print("Versioning models...")
    
    version_id = f"v{datetime.now().isoformat().replace(':', '')}"
    
    # Save to manifest
    manifest = {
        "active_version": version_id,
        "versions": {
            version_id: {
                "metrics": results,
                "created_at": datetime.now().isoformat()
            }
        }
    }
    
    with open("models/manifest.json", "w") as f:
        json.dump(manifest, f, indent=2)
    
    print(f"✓ Versioned models: {version_id}")
    return version_id

@flow(name="bitcoin-ml-pipeline")
def ml_training_flow(data_path: str = "data/raw/bitcoin_timeseries.csv"):
    """
    Complete ML training pipeline with orchestration
    
    Steps:
    1. Fetch Bitcoin data
    2. Validate data quality
    3. Engineer features
    4. Preprocess data
    5. Split into train/test
    6. Train all models (in parallel)
    7. Evaluate performance
    8. Version best models
    """
    
    print("\n" + "="*60)
    print("🚀 BITCOIN ML TRAINING PIPELINE")
    print("="*60 + "\n")
    
    # Data pipeline
    df_raw = fetch_data_task(days=365)
    is_valid = validate_data_task(df_raw)
    
    if not is_valid:
        raise ValueError("Data validation failed")
    
    df_features = engineer_features_task(df_raw)
    df_processed, scaler = preprocess_task(df_features)
    X_train, X_test = split_data_task(df_processed, test_size=0.2)
    
    # Create targets
    y_train_clf = (df_processed['price'].shift(-1) > df_processed['price']).astype(int)
    y_train_reg = df_processed['price'].pct_change()
    
    # Training pipeline
    results = train_models_task(X_train, X_test, y_train_clf, y_train_reg)
    evaluated = evaluate_models_task(results)
    version_id = version_models_task(evaluated, models={})
    
    print("\n" + "="*60)
    print(f"✅ PIPELINE COMPLETE - Models versioned as {version_id}")
    print("="*60 + "\n")
    
    return {"version": version_id, "metrics": evaluated}

if __name__ == "__main__":
    # Execute flow
    state = ml_training_flow(data_path="data/raw/bitcoin_timeseries.csv")
    print(state)
\end{lstlisting}

\section{Task Specifications}

\subsection{Task 1: Fetch Data}

\begin{table}[H]
\centering
\caption{Task: Fetch Data}
\begin{tabularx}{\textwidth}{lX}
\toprule
\textbf{Parameter} & \textbf{Value} \\
\midrule
Name & fetch\_data\_task \\
Retries & 3 (exponential backoff) \\
Retry Delay & 30 seconds \\
Timeout & 300 seconds (5 minutes) \\
Input & days (int, default=365) \\
Output & pd.DataFrame (OHLCV data) \\
\bottomrule
\end{tabularx}
\end{table}

\subsection{Task 2: Validate Data}

\begin{enumerate}
    \item Check non-empty dataframe
    \item Verify no missing price values
    \item Validate price ranges (> 0)
    \item Check date ordering
\end{enumerate}

\subsection{Task 3: Engineer Features}

\begin{enumerate}
    \item Calculate 24 technical indicators
    \item Compute temporal features
    \item Drop invalid rows (NaN)
    \item Return feature matrix
\end{enumerate}

\subsection{Task 4: Preprocess}

\begin{enumerate}
    \item Standardize features (StandardScaler)
    \item Handle missing values
    \item Return (processed data, scaler)
\end{enumerate}

\subsection{Task 5: Split Data}

\begin{enumerate}
    \item Temporal split (preserving order)
    \item Default 80/20 ratio
    \item Return (X\_train, X\_test)
\end{enumerate}

\subsection{Task 6: Train Models}

\begin{itemize}
    \item \textbf{Parallel:} Classification, Regression, Deep Learning
    \item \textbf{Retries:} 2 (OOM recovery)
    \item \textbf{Timeout:} 1800 seconds (30 minutes)
    \item \textbf{Returns:} Dict of results
\end{itemize}

\subsection{Task 7: Evaluate}

\begin{enumerate}
    \item Compute metrics on test set
    \item Compare with baseline
    \item Detect data drift
    \item Return evaluation results
\end{enumerate}

\subsection{Task 8: Version Models}

\begin{enumerate}
    \item Create version ID (timestamp)
    \item Save to manifest.json
    \item Upload model artifacts
    \item Return version string
\end{enumerate}

\section{Error Handling}

\subsection{Retry Strategy}

\begin{lstlisting}[language=Python, caption={Retry Configuration}]
@task(
    retries=3,                      # Retry up to 3 times
    retry_delay_seconds=30,         # Wait 30s between retries
    retry_jitter=True,              # Add randomness (0-10s)
    timeout_seconds=300             # Fail after 5 minutes
)
def fetch_data_task(days: int = 365):
    try:
        return fetch_bitcoin_data(days=days)
    except requests.exceptions.Timeout:
        raise  # Will trigger retry
    except requests.exceptions.ConnectionError:
        raise  # Will trigger retry
    except Exception as e:
        # Log unexpected errors
        logger.error(f"Unexpected error: {e}")
        raise  # Re-raise for Prefect to handle
\end{lstlisting}

\subsection{Exception Handling}

\begin{table}[H]
\centering
\caption{Exception Handling Strategy}
\begin{tabularx}{\textwidth}{lXc}
\toprule
\textbf{Exception} & \textbf{Action} & \textbf{Retries} \\
\midrule
Connection Error & Retry with backoff & 3 \\
Timeout & Retry with delay & 3 \\
OOM Error & Retry with smaller batch & 2 \\
Data Validation Error & Fail immediately & 0 \\
Model Training Error & Retry once & 1 \\
\bottomrule
\end{tabularx}
\end{table}

\section{Deployment Configuration}

\subsection{Deployment Manifest}

\begin{lstlisting}[language=yaml, caption={Prefect Deployment Config}]
deployments:
  - name: bitcoin-ml-daily
    flow_name: bitcoin-ml-pipeline
    schedule:
      cron: '0 2 * * *'  # Daily at 2 AM UTC
    parameters:
      data_path: data/raw/bitcoin_timeseries.csv
    work_queue: default
    tags:
      - production
      - daily
    description: Daily model retraining with latest Bitcoin data
    
  - name: bitcoin-ml-manual
    flow_name: bitcoin-ml-pipeline
    schedule: null  # Manual trigger only
    parameters:
      data_path: data/raw/bitcoin_timeseries.csv
    work_queue: default
    tags:
      - manual
    description: Manual model training execution
\end{lstlisting}

\subsection{Execution Modes}

\subsubsection{Mode 1: Local Execution}
\begin{lstlisting}[language=bash]
python prefect/flows/ml_pipeline.py
\end{lstlisting}

\subsubsection{Mode 2: Prefect Server}
\begin{lstlisting}[language=bash]
# Start Prefect server
prefect server start

# Deploy flow
prefect deployment build prefect/flows/ml_pipeline.py:ml_training_flow -n bitcoin-ml-daily

# Execute via server
prefect deployment run bitcoin-ml-pipeline/bitcoin-ml-daily
\end{lstlisting}

\subsubsection{Mode 3: Prefect Cloud}
\begin{lstlisting}[language=bash]
# Authenticate with Prefect Cloud
prefect cloud login

# Deploy to cloud
prefect deployment push prefect/flows/ml_pipeline.py:ml_training_flow -n bitcoin-ml-daily

# Trigger from cloud
prefect deployment run bitcoin-ml-pipeline/bitcoin-ml-daily
\end{lstlisting}

\section{Monitoring \& Observability}

\subsection{Flow Runs Dashboard}

Prefect Server provides a web UI showing:
\begin{itemize}
    \item Flow run history
    \item Task execution status
    \item Execution logs
    \item Performance metrics
\end{itemize}

\subsection{Logging}

\begin{lstlisting}[language=Python, caption={Structured Logging}]
from prefect import get_run_logger

@task
def my_task():
    logger = get_run_logger()
    logger.info("Task started")
    logger.debug(f"Processing {n} records")
    logger.warning("Potential issue detected")
    logger.error("Error occurred")
\end{lstlisting}

\subsection{Notifications}

Prefect can send notifications on:
\begin{itemize}
    \item Flow completion
    \item Task failures
    \item SLA violations
    \item Custom events
\end{itemize}

\section{Performance Characteristics}

\begin{table}[H]
\centering
\caption{Pipeline Performance}
\begin{tabularx}{\textwidth}{lccc}
\toprule
\textbf{Stage} & \textbf{Duration} & \textbf{CPU} & \textbf{Memory} \\
\midrule
Fetch Data & 2-3 min & Low & Low \\
Validate Data & 30 sec & Low & Low \\
Engineer Features & 1-2 min & Medium & Medium \\
Preprocess & 1-2 min & Medium & Medium \\
Split Data & 30 sec & Low & Low \\
Train Models & 10-15 min & High & High \\
Evaluate & 2-3 min & Medium & Medium \\
Version & 1 min & Low & Low \\
\midrule
\textbf{Total} & \textbf{18-30 min} & \textbf{} & \textbf{} \\
\bottomrule
\end{tabularx}
\end{table}

\section{Summary}

Prefect provides:
\begin{itemize}
    \item ✅ Task orchestration with dependencies
    \item ✅ Automatic retry logic with exponential backoff
    \item ✅ Comprehensive logging and monitoring
    \item ✅ Scheduled and manual execution
    \item ✅ Error handling and recovery
    \item ✅ Cloud-ready infrastructure
\end{itemize}

